% !TeX root = ../main.tex

\chapter{Chosen Conceptual Design}
\label{chap:con-design}
The major system functions are summarized below.

\begin{enumerate}
	\item \textbf{Power-Supply Interface:} Connect to the six monitored rails and
	route each rail through the associated voltage/current monitoring channel.
	\item \textbf{Current Measurement:} Convert load current to a differential
	voltage across a precision shunt resistor and digitize it with the
	LTC2945.
	\item \textbf{Voltage Measurement:} Measure the monitored rail voltage using
	the LTC2945 supply-voltage ADC channel.
	\item \textbf{Electrical Isolation:} Maintain galvanic isolation between each
	floating monitored supply and the common controller-side I\textsuperscript{2}C bus.
	\item \textbf{Digital Acquisition:} Poll each monitor, identify communication
	faults, apply channel-specific sign/scaling information, and calculate
	engineering-unit values.
	\item \textbf{Alarm Processing:} Compare voltage and current values against
	configured upper and lower limits and generate channel alarm states.
	\item \textbf{Local User Interface:} Present all rail measurements, alarms,
	network state, and communication status on the Arduino GIGA Display.
	\item \textbf{Network/EPICS Interface:} Publish measurements, alarm limits,
	alarm state, and connection status over Ethernet for the EPICS IOC and
	Phoebus interface.
\end{enumerate}

\subsection{Functional Block Diagram}

\begin{figure}[H]
	\centering
	
	\makebox[\textwidth][c]{%
		\rotatebox{90}{%
			\begin{minipage}{0.9\textheight}
				\centering
				
				\resizebox{0.78\textheight}{!}{%
					\begin{tikzpicture}[
						>=Stealth,
						font=\small,
						block/.style={
							draw,
							rectangle,
							rounded corners=2pt,
							align=center,
							minimum height=1.25cm,
							minimum width=2.75cm,
							text width=2.45cm,
							inner sep=5pt,
							line width=0.8pt
						},
						rail/.style={
							draw,
							rectangle,
							rounded corners=2pt,
							align=center,
							minimum height=0.55cm,
							minimum width=1.35cm,
							font=\footnotesize,
							line width=0.7pt
						},
						group/.style={
							draw,
							dashed,
							rounded corners=4pt,
							inner sep=8pt,
							line width=0.8pt
						},
						signal/.style={
							->,
							line width=0.9pt
						},
						line/.style={
							line width=0.9pt
						},
						bus/.style={
							line width=1.1pt
						},
						busarrow/.style={
							->,
							line width=1.1pt
						},
						siglabel/.style={
							font=\scriptsize,
							fill=white,
							inner sep=1.5pt,
							text depth=0pt,
							text height=1.5ex,
							align=center
						}
						]
						
						% ============================================================
						% MONITORED POWER RAILS
						% ============================================================
						
						\node[rail] (p6a) at (0, 2.25) {+6 V A};
						\node[rail] (p6b) at (0, 1.35) {+6 V B};
						\node[rail] (p5)  at (0, 0.45) {+5 V};
						\node[rail] (m5)  at (0,-0.45) {$-5$ V};
						\node[rail] (p15) at (0,-1.35) {+15 V};
						\node[rail] (m15) at (0,-2.25) {$-15$ V};
						
						\node[
						above=0.30cm of p6a,
						font=\bfseries
						] {Monitored Power-Supply Rails};
						
						
						% ============================================================
						% RAIL COLLECTION BUS
						% ============================================================
						
						\coordinate (railBusTop)    at (1.15, 2.25);
						\coordinate (railBusBottom) at (1.15,-2.25);
						\coordinate (railBusMid)    at (1.15, 0);
						
						\draw[bus] (railBusTop) -- (railBusBottom);
						
						% Individual rail connections -- no overlapping arrowheads.
						\foreach \railname/\y in {
							p6a/2.25,
							p6b/1.35,
							p5/0.45,
							m5/-0.45,
							p15/-1.35,
							m15/-2.25
						}{
							\draw[line] (\railname.east) -- (1.15,\y);
						}
						
						
						% ============================================================
						% POWER-SUPPLY INTERFACE BOARD
						% ============================================================
						
						\node[
						block,
						anchor=west
						] (shunts) at (2.15,0)
						{
							High-Side\\
							Current-Sense Shunts\\[2pt]
							\footnotesize
							One independent\\
							sense element per rail
						};
						
						\node[
						block,
						right=0.65cm of shunts
						] (ltc)
						{
							LTC2945\\
							Voltage / Current\\
							Monitor Channels\\[2pt]
							\footnotesize
							Six independent\\
							channels
						};
						
						\node[
						block,
						right=0.65cm of ltc
						] (iso)
						{
							I\textsuperscript{2}C Digital\\
							Isolation\\[2pt]
							\footnotesize
							One isolated interface\\
							per monitored rail
						};
						
						
						\draw[busarrow]
						(railBusMid)
						--
						(shunts.west);
						
						
						\draw[busarrow]
						(shunts.east)
						--
						node[
						siglabel,
						midway,
						rotate=90,
						above=6pt
						]
						{Sense Voltage}
						(ltc.west);
						
						
						\draw[busarrow]
						(ltc.east)
						--
						node[
						siglabel,
						midway,
						above=4pt
						]
						{I\textsuperscript{2}C}
						(iso.west);
						
						
						% ============================================================
						% POWER-SUPPLY INTERFACE GROUP
						% ============================================================
						
						\begin{scope}[on background layer]
							\node[
							group,
							fit=(shunts)(ltc)(iso)
							] (interfaceGroup) {};
						\end{scope}
						
						\node[
						font=\bfseries,
						fill=white,
						inner sep=2pt,
						anchor=south
						]
						at (interfaceGroup.north)
						{Power Supply Interface Board};
						
						
						% ============================================================
						% ISOLATION BARRIER
						% ============================================================
						
						\coordinate (barrierTop)
						at ($(iso.north east)+(0.50,0.55)$);
						
						\coordinate (barrierBottom)
						at ($(iso.south east)+(0.50,-0.55)$);
						
						% Keep only the dashed isolation-barrier line; no text label.
						\draw[
						dashed,
						line width=1.1pt
						]
						(barrierTop) -- (barrierBottom);
						
						
						% ============================================================
						% ARDUINO GIGA
						% ============================================================
						
						\node[
						block,
						right=1.45cm of iso
						] (giga)
						{
							Arduino GIGA R1 WiFi\\[2pt]
							\footnotesize
							Acquire Measurements\\
							Apply Scaling and Polarity\\
							Evaluate Alarm Limits\\
							Manage Configuration
						};
						
						
						\draw[busarrow]
						(iso.east)
						--
						node[
						siglabel,
						midway,
						rotate=90,
						above=6pt
						]
						{Isolated I\textsuperscript{2}C}
						(giga.west);
						
						
						% ============================================================
						% DISPLAY AND ETHERNET
						% ============================================================
						
						\node[
						block,
						right=1.25cm of giga,
						yshift=1.55cm
						] (display)
						{
							GIGA Display Shield\\[2pt]
							\footnotesize
							Voltage / Current\\
							Alarm Status\\
							Communication Status
						};
						
						\node[
						block,
						right=1.25cm of giga,
						yshift=-1.55cm
						] (ethernet)
						{
							Ethernet Interface\\[2pt]
							\footnotesize
							PSC Protocol\\
							Configuration Interface
						};
						
						\coordinate (controllerBranch)
						at ($(giga.east)+(0.55,0)$);
						
						\draw[line]
						(giga.east)
						--
						(controllerBranch);
						
						\draw[signal]
						(controllerBranch)
						|-
						(display.west);
						
						\draw[signal]
						(controllerBranch)
						|-
						(ethernet.west);
						
						\fill (controllerBranch) circle (1.5pt);
						
						
						% ============================================================
						% CONTROLLER GROUP
						% ============================================================
						
						\begin{scope}[on background layer]
							\node[
							group,
							fit=(giga)(display)(ethernet)
							] (controllerGroup) {};
						\end{scope}
						
						\node[
						font=\bfseries,
						fill=white,
						inner sep=2pt,
						anchor=south
						]
						at (controllerGroup.north)
						{Controller Assembly};
						
						
						% ============================================================
						% EPICS IOC
						% ============================================================
						
						\node[
						block,
						right=1.20cm of ethernet
						] (ioc)
						{
							EPICS IOC / PSCDriver\\[2pt]
							\footnotesize
							Receive Measurements\\
							Publish Engineering Values\\
							Generate EPICS Alarms
						};
						
						\draw[busarrow]
						(ethernet.east)
						--
						node[
						siglabel,
						midway,
						rotate=90,
						above=6pt
						]
						{Ethernet}
						(ioc.west);
						
						
						% ============================================================
						% PHOEBUS
						% ============================================================
						
						\node[
						block,
						right=1.20cm of ioc
						] (phoebus)
						{
							CS-Studio Phoebus\\[2pt]
							\footnotesize
							Operator Interface\\
							Voltage / Current Display\\
							Alarm and System Status
						};
						
						\draw[busarrow]
						(ioc.east)
						--
						node[
						siglabel,
						midway,
						rotate=90,
						below=6pt
						]
						{EPICS Channel Access}
						(phoebus.west);
						
						
						% ============================================================
						% SUPERVISORY GROUP
						% ============================================================
						
						\begin{scope}[on background layer]
							\node[
							group,
							fit=(ioc)(phoebus)
							] (supervisoryGroup) {};
						\end{scope}
						
						\node[
						font=\bfseries,
						fill=white,
						inner sep=2pt,
						anchor=south
						]
						at (supervisoryGroup.north)
						{Supervisory Control System};
						
					\end{tikzpicture}%
				}
				
				\vspace{0.25cm}
				
				\caption{
					Functional block diagram of the multi-rail isolated
					power supply monitoring system. Six independent power-supply rails
					are measured using high-side current sensing and LTC2945 voltage/current
					monitoring channels. Digital isolation preserves electrical independence
					between the monitored supplies and the common controller. The Arduino
					GIGA processes the measurements, provides a local operator display, and
					transmits measurement and status information over Ethernet to the EPICS
					IOC and CS-Studio Phoebus operator interface.
				}
				\label{fig:functional-block-diagram}
				
			\end{minipage}%
		}%
	}
	
\end{figure}